\documentclass[a4paper,12pt]{article}
\usepackage[utf8]{inputenc}
\usepackage[T1]{fontenc}
\usepackage[italian]{babel}
\usepackage{amsmath}
\usepackage{amssymb}
\usepackage{amsthm}
\usepackage{graphicx}
\usepackage{geometry}
\usepackage{hyperref}
\usepackage{booktabs}

\geometry{top=2.5cm, bottom=2.5cm, left=3cm, right=3cm}

\theoremstyle{definition}
\newtheorem{definizione}{Definizione}[section]
\newtheorem{osservazione}{Osservazione}[section]

\theoremstyle{remark}
\newtheorem*{intuizione}{Intuizione geometrica/numerica}

\title{\textbf{Approssimazione Numerica di Integrali}\\[0.5em]
\large Guida di Studio ---
Metodi Matematici e Numerici per Ingegneria Informatica}
\author{Fausto Vanella}
\date{Giugno 2026}

\begin{document}

\maketitle
\tableofcontents
\newpage

% =============================================================================
\section{Formula di Quadratura}
% =============================================================================

\subsection{Definizione e motivazione}

Quando si parla di \textbf{formula di quadratura}, ci si riferisce
all'approssimazione numerica di un integrale definito.
Considerato l'integrale di una funzione $f(x)$ sull'intervallo $[a,b]$,
è possibile approssimarlo come somma pesata dei valori di $f$ in un insieme
finito di punti, più un termine di errore residuo:

\begin{equation}
  \int_a^b f(x)\,dx = \sum_{i=0}^{n} \omega_i\,f(x_i) + R_n
  \label{eq:quadratura}
\end{equation}

\noindent dove:
\begin{itemize}
  \item $\{x_i\}_{i=0}^{n}$ sono gli $n+1$ \textbf{nodi} di quadratura:
        i punti dell'intervallo $[a,b]$ nei quali la funzione $f$ viene valutata;
  \item $\{\omega_i\}_{i=0}^{n}$ sono i \textbf{pesi} di quadratura:
        coefficienti numerici che determinano quanto ciascun valore $f(x_i)$
        contribuisce all'approssimazione;
  \item $R_n$ è il \textbf{resto} (o errore di quadratura):
        lo scarto tra il valore esatto dell'integrale e l'approssimazione.
\end{itemize}

\begin{definizione}[Formula di Quadratura Numerica]
Una qualunque formula della forma~\eqref{eq:quadratura} si dice di
\textbf{quadratura numerica}: essa stima l'integrale definito di $f(x)$
sull'intervallo $[a,b]$ senza richiedere il calcolo esplicito della
primitiva di $f$.
\end{definizione}

\subsection{Grado di precisione}

\begin{definizione}[Grado di Precisione]
Il \textbf{grado di precisione} $p$ di una formula di quadratura è la
misura della sua accuratezza. In particolare, si dice che la formula ha
grado di precisione $p$ se:
\begin{enumerate}
  \item applicandola a qualsiasi polinomio di grado $\le p$ fornisce il
        risultato esatto (ovvero $R_n = 0$);
  \item esiste almeno un polinomio di grado $p+1$ per cui la formula
        \emph{non} è esatta.
\end{enumerate}
\end{definizione}

\subsection{Formule di tipo interpolatorio}

Le formule di quadratura che studieremo sono di tipo \textbf{interpolatorio}.
L'idea fondamentale è: anziché integrare direttamente $f(x)$,
la si sostituisce con un polinomio $P_n(x)$ di grado $\le n$ che passa
per gli $n+1$ nodi, e si integra quest'ultimo.
L'errore commesso è l'integrale del termine di errore dell'interpolazione
polinomiale $E_n(x)$.

Formalmente, dato che $f(x) = P_n(x) + E_n(x)$ per costruzione, possiamo
scrivere:
\begin{equation}
  \int_a^b f(x)\,dx
  = \underbrace{\int_a^b P_n(x)\,dx}_{\text{parte approssimata}}
  + \underbrace{\int_a^b E_n(x)\,dx}_{R_n}
  \label{eq:decomposizione}
\end{equation}

\noindent dove $P_n(x)$ è il polinomio interpolante passante per gli
$n+1$ nodi $x_0, x_1, \ldots, x_n$.

La formula risulta \textbf{esatta} se e solo se il resto si azzera:
\[
  R_n = \int_a^b E_n(x)\,dx = 0
\]

% =============================================================================
\section{Newton-Cotes: Nodi Equidistanti}
% =============================================================================

\subsection{Impostazione}

Le \textbf{formule di Newton-Cotes} sono formule interpolatorie in cui i
nodi sono \textbf{equidistanti}, con il primo nodo coincidente con
l'estremo sinistro $a$ e l'ultimo con l'estremo destro $b$:
\[
  x_0 = a, \quad x_n = b, \quad
  x_i = x_0 + i\,h \quad \text{per } i = 0,1,\ldots,n
\]

\noindent dove $h = \dfrac{b-a}{n}$ è il \textbf{passo di discretizzazione},
cioè la distanza uniforme tra nodi consecutivi.

\begin{intuizione}
Suddividere uniformemente $[a,b]$ semplifica i calcoli: la distanza tra
qualsiasi coppia di nodi $x_i$ e $x_j$ è un multiplo intero di $h$.
Questo permette di esprimere tutte le differenze in forma compatta e di
ricondurre i polinomi di Lagrange (e i relativi integrali) a espressioni
polinomiali in una variabile adimensionale $t$.
\end{intuizione}

\subsection{Formule per le differenze tra nodi}

Prendendo un punto generico $x = x_0 + t h$
(dove $t$ è un parametro reale), calcoliamo la differenza tra $x$ e il
generico nodo $x_i = x_0 + i\,h$:
\begin{equation}
  x - x_i = (x_0 + th) - (x_0 + ih) = h(t-i)
  \label{eq:diff-x-xi}
\end{equation}

Analogamente, la differenza tra due nodi $x_i$ e $x_j$ è:
\begin{equation}
  x_i - x_j = (x_0 + ih) - (x_0 + jh) = h(i-j)
  \label{eq:diff-xi-xj}
\end{equation}

Queste due relazioni sono fondamentali: ci permetteranno, nella Sezione~3,
di cancellare sistematicamente i fattori $h$ nei polinomi di Lagrange e di
ricavare un'espressione chiusa per i pesi $\omega_i$.

% =============================================================================
\section{Derivazione dei Pesi tramite Polinomi di Lagrange}
% =============================================================================

\subsection{Integrale come combinazione di integrali di Lagrange}

Supponendo che il polinomio approssimante $P_n(x)$ nella
decomposizione~\eqref{eq:decomposizione} sia il polinomio di Lagrange
passante per gli $n+1$ nodi:
\begin{equation}
  P_n(x) = \sum_{i=0}^{n} l_i(x)\,f(x_i)
  \label{eq:lagrange-pn}
\end{equation}

\noindent dove i \textbf{polinomi fondamentali di Lagrange} $l_i(x)$ sono
definiti come:
\begin{equation}
  l_i(x) = \frac{(x-x_0)\cdots(x-x_{i-1})(x-x_{i+1})\cdots(x-x_n)}
                 {(x_i-x_0)\cdots(x_i-x_{i-1})(x_i-x_{i+1})\cdots(x_i-x_n)}
  \label{eq:lagrange-li}
\end{equation}

\noindent e soddisfano la proprietà $l_i(x_j) = \delta_{ij}$
(cioè: $l_i(x_i) = 1$ e $l_i(x_j) = 0$ per $j \ne i$).

Sostituendo~\eqref{eq:lagrange-pn} nell'integrale e sfruttando la
linearità dell'integrazione (ovvero lo scambio tra $\int$ e $\sum$):
\[
  \int_a^b f(x)\,dx
  \approx \int_a^b \sum_{i=0}^{n} l_i(x)\,f(x_i)\,dx
  = \sum_{i=0}^{n} f(x_i)\int_a^b l_i(x)\,dx
\]

\noindent Poiché $f(x_i)$ è un valore costante rispetto alla variabile di
integrazione $x$, può uscire dall'integrale.
Ne consegue direttamente che i \textbf{pesi di quadratura} sono:
\begin{equation}
  \boxed{\omega_i = \int_a^b l_i(x)\,dx = \int_{x_0}^{x_n} l_i(x)\,dx}
  \label{eq:pesi}
\end{equation}

I pesi dipendono solo dalla scelta dei nodi, non dalla funzione $f$:
questo è il vantaggio delle formule interpolatorie.

\subsection{Cambio di variabile per il calcolo esplicito dei pesi}

Per calcolare analiticamente i pesi~\eqref{eq:pesi}, effettuiamo il
cambio di variabile:
\[
  x = x_0 + t\,h \;\Longrightarrow\; dx = h\,dt
\]

Gli estremi di integrazione si trasformano come segue:
\begin{align*}
  x = a = x_0 &\;\Rightarrow\;
    t = \frac{x - x_0}{h} = \frac{x_0 - x_0}{h} = 0 \\[4pt]
  x = b = x_n &\;\Rightarrow\;
    t = \frac{x - x_0}{h} = \frac{x_n - x_0}{h} = \frac{n\,h}{h} = n
\end{align*}

\noindent Ora, usando le relazioni~\eqref{eq:diff-x-xi}
e~\eqref{eq:diff-xi-xj}, riscriviamo i polinomi di Lagrange in funzione
di $t$. Nel prodotto~\eqref{eq:lagrange-li}, ogni fattore diventa:
\[
  \frac{x - x_j}{x_i - x_j} = \frac{h(t-j)}{h(i-j)} = \frac{t-j}{i-j}
\]

\noindent I fattori $h$ al numeratore e al denominatore si cancellano
esattamente. Quindi:
\[
  l_i(x) = \prod_{\substack{j=0\\j\ne i}}^{n} \frac{t-j}{i-j}
\]

Sostituendo nell'espressione del peso:
\begin{align*}
  \omega_i
  &= \int_{x_0}^{x_n} l_i(x)\,dx
   = h \int_0^n \prod_{\substack{j=0\\j\ne i}}^{n} \frac{t-j}{i-j}\,dt
\end{align*}

\noindent Poiché i termini $(i-j)$ al denominatore sono costanti rispetto
a $t$, si raccolgono fuori dall'integrale:

\begin{equation}
  \boxed{
    \omega_i = h \int_0^n \prod_{\substack{j=0\\j\ne i}}^{n}
    \frac{t-j}{i-j}\,dt
  }
  \label{eq:pesi-t}
\end{equation}

\noindent dove:
\begin{itemize}
  \item $h = (b-a)/n$ è il passo di discretizzazione;
  \item il prodotto esclude il termine $j = i$ (che darebbe divisione per
        zero al denominatore);
  \item l'integrale è calcolato nella variabile $t \in [0, n]$,
        corrispondente a $x \in [x_0, x_n] = [a,b]$.
\end{itemize}

% =============================================================================
\section{Newton-Cotes Semplici}
% =============================================================================

\subsection{Formula del Trapezio (Newton-Cotes a 2 Nodi)}

La \textbf{formula del trapezio} corrisponde a Newton-Cotes con $n = 1$:
si utilizzano soltanto 2 nodi, i due estremi dell'intervallo:
\[
  x_0 = a, \quad x_1 = b, \quad h = b - a
\]

\subsubsection*{Calcolo dei pesi}

Con $n = 1$, la formula~\eqref{eq:pesi-t} richiede di integrare in
$t \in [0, 1]$.

\medskip\noindent
\textbf{Peso $\omega_0$} (associato al nodo $x_0 = a$, dove $l_0$
vale 1 e si azzera in $x_1$):
\begin{align*}
  \omega_0
  &= h \int_0^1 \frac{t - 1}{0 - 1}\,dt
   = h \int_0^1 \frac{t-1}{-1}\,dt
   = h \int_0^1 (1-t)\,dt \\
  &= h \left[t - \frac{t^2}{2}\right]_0^1
   = h \left(1 - \frac{1}{2}\right)
   = \frac{h}{2}
\end{align*}

\noindent
\textbf{Peso $\omega_1$} (associato al nodo $x_1 = b$):
\begin{align*}
  \omega_1
  &= h \int_0^1 \frac{t - 0}{1 - 0}\,dt
   = h \int_0^1 t\,dt
   = h \left[\frac{t^2}{2}\right]_0^1
   = \frac{h}{2}
\end{align*}

\noindent I due pesi sono uguali tra loro e coincidono con la metà del
passo: ciascun estremo contribuisce in egual misura all'approssimazione.

\subsubsection*{Formula}

\begin{equation}
  \boxed{
    \int_a^b f(x)\,dx
    \approx \frac{h}{2}\bigl[f(x_0) + f(x_1)\bigr]
    = \frac{b-a}{2}\bigl[f(a) + f(b)\bigr]
  }
  \label{eq:trapezio}
\end{equation}

\begin{intuizione}
Questa formula coincide esattamente con la formula per l'area di un
trapezio ``capovolto'': $f(a)$ e $f(b)$ sono le due basi parallele,
$h = b - a$ è l'altezza.
Si approssima quindi l'area sottesa alla curva $f(x)$ con l'area del
trapezio delimitato dalla retta secante passante per i due punti
estremi $(a,f(a))$ e $(b,f(b))$.
\end{intuizione}

\subsubsection*{Errore della formula del trapezio}

Il termine di errore dell'interpolazione polinomiale con $n+1$ nodi vale:
\[
  E_n(x) = \prod_{i=0}^{n}(x - x_i)\,\frac{f^{(n+1)}(\xi)}{(n+1)!}
\]

\noindent dove $\xi = \xi(x)$ è un punto (incognito) appartenente
all'intervallo $[a,b]$.
Integrando su $[a,b]$ (e assumendo che il segno di
$f^{(n+1)}$ non cambi, per poter fattorizzare), il resto è:
\begin{equation}
  R_n = \int_a^b E_n(x)\,dx
      = \frac{f^{(n+1)}(\xi)}{(n+1)!}
        \int_a^b (x-x_0)(x-x_1)\cdots(x-x_n)\,dx
  \label{eq:errore-gen}
\end{equation}

\noindent Per $n = 1$ (due nodi $x_0 = a$, $x_1 = b$):
\[
  R_1 = \frac{f''(\xi)}{2!}\int_a^b (x - a)(x - b)\,dx
      = \frac{f''(\xi)}{2}\int_a^b (x-a)(x-b)\,dx
\]

\noindent Calcoliamo l'integrale con la sostituzione $u = x - a$
(da cui $x - b = u - h$, $du = dx$, estremi $0 \to h$):
\begin{align*}
  \int_a^b (x-a)(x-b)\,dx
  &= \int_0^h u(u - h)\,du
   = \int_0^h (u^2 - hu)\,du \\
  &= \left[\frac{u^3}{3} - \frac{h\,u^2}{2}\right]_0^h
   = \frac{h^3}{3} - \frac{h^3}{2}
   = h^3\!\left(\frac{2-3}{6}\right)
   = -\frac{h^3}{6}
\end{align*}

\noindent Quindi:
\[
  R_1 = \frac{f''(\xi)}{2} \cdot \left(-\frac{h^3}{6}\right)
\]

\begin{equation}
  \boxed{R_1 = -\frac{h^3}{12}\,f''(\xi)}
  \label{eq:errore-trapezio}
\end{equation}

\noindent dove $\xi \in (a,b)$.
Poiché il segno del resto è opposto a quello di $f''(\xi)$:
\begin{itemize}
  \item se $f$ è \textbf{convessa} ($f'' > 0$): la formula del trapezio
        \emph{sovrastima} l'integrale reale;
  \item se $f$ è \textbf{concava} ($f'' < 0$): la formula del trapezio
        \emph{sottostima} l'integrale reale.
\end{itemize}

Il \textbf{grado di precisione} della formula del trapezio è $p = 1$:
la formula è esatta per tutte le funzioni polinomiali di grado $\le 1$
(funzioni costanti e rette), ma non per il polinomio $x^2$.

\subsection{Formula di Cavalieri-Simpson (Newton-Cotes a 3 Nodi)}

La \textbf{formula di Cavalieri-Simpson} corrisponde a Newton-Cotes con
$n = 2$: si usano 3 nodi equidistanti, compresi gli estremi e il punto
medio dell'intervallo:
\[
  x_0 = a, \quad
  x_1 = \frac{a+b}{2}, \quad
  x_2 = b, \quad
  h = \frac{b-a}{2}
\]

\noindent dove $h$ è ora la metà dell'ampiezza dell'intervallo.

\subsubsection*{Calcolo dei pesi}

Con $n = 2$, la formula~\eqref{eq:pesi-t} richiede di integrare in
$t \in [0, 2]$.

\medskip\noindent
\textbf{Peso $\omega_0$} (associato a $x_0$):
\begin{align*}
  \omega_0
  &= h\int_0^2 \frac{(t-1)(t-2)}{(0-1)(0-2)}\,dt
   = h\int_0^2 \frac{(t-1)(t-2)}{(-1)(-2)}\,dt
   = \frac{h}{2}\int_0^2 (t-1)(t-2)\,dt \\
  &= \frac{h}{2}\int_0^2 (t^2 - 3t + 2)\,dt
   = \frac{h}{2}\left[\frac{t^3}{3} - \frac{3t^2}{2} + 2t\right]_0^2 \\
  &= \frac{h}{2}\!\left(\frac{8}{3} - 6 + 4\right)
   = \frac{h}{2}\!\left(\frac{8}{3} - 2\right)
   = \frac{h}{2} \cdot \frac{8 - 6}{3}
   = \frac{h}{2} \cdot \frac{2}{3}
   = \frac{h}{3}
\end{align*}

\noindent
\textbf{Peso $\omega_1$} (associato a $x_1$, il punto medio):
\begin{align*}
  \omega_1
  &= h\int_0^2 \frac{(t-0)(t-2)}{(1-0)(1-2)}\,dt
   = h\int_0^2 \frac{t(t-2)}{1 \cdot (-1)}\,dt
   = -h\int_0^2 (t^2 - 2t)\,dt \\
  &= -h\left[\frac{t^3}{3} - t^2\right]_0^2
   = -h\!\left(\frac{8}{3} - 4\right)
   = -h \cdot \frac{8 - 12}{3}
   = -h \cdot \left(-\frac{4}{3}\right)
   = \frac{4h}{3}
\end{align*}

\noindent
\textbf{Peso $\omega_2$} (associato a $x_2$):
\begin{align*}
  \omega_2
  &= h\int_0^2 \frac{(t-0)(t-1)}{(2-0)(2-1)}\,dt
   = \frac{h}{2}\int_0^2 t(t-1)\,dt
   = \frac{h}{2}\int_0^2 (t^2 - t)\,dt \\
  &= \frac{h}{2}\left[\frac{t^3}{3} - \frac{t^2}{2}\right]_0^2
   = \frac{h}{2}\!\left(\frac{8}{3} - 2\right)
   = \frac{h}{2} \cdot \frac{2}{3}
   = \frac{h}{3}
\end{align*}

\noindent I pesi risultanti sono $(\omega_0,\omega_1,\omega_2) =
(h/3,\; 4h/3,\; h/3)$.
Il nodo centrale $x_1$ riceve un peso quattro volte maggiore degli
estremi: questo riflette il fatto che, per catturare la curvatura di
$f$, il valore al centro è più ``informativo'' dei valori agli estremi.

\subsubsection*{Formula}

\begin{equation}
  \boxed{
    \int_a^b f(x)\,dx
    \approx \frac{h}{3}\bigl[f(x_0) + 4f(x_1) + f(x_2)\bigr]
    = \frac{b-a}{6}\left[f(a) + 4f\!\left(\tfrac{a+b}{2}\right) + f(b)\right]
  }
  \label{eq:simpson}
\end{equation}

\begin{intuizione}
Rispetto al trapezio, qui $f$ viene approssimata con una \emph{parabola}
(polinomio di grado 2) passante per i tre punti
$(x_0,f(x_0))$, $(x_1,f(x_1))$, $(x_2,f(x_2))$.
Il peso quadruplo al centro nasce dall'integrale di $l_1(t)$, che è
una parabola con vertice ampio: la parabola interpolante ``piega'' verso
il punto medio catturando la curvatura di $f$ in modo molto più fedele
di una retta.
\end{intuizione}

% =============================================================================
\section{Formule di Quadratura Composte}
% =============================================================================

\subsection{Motivazione}

Le formule semplici viste finora (trapezio, Cavalieri-Simpson) diventano
inaccurate quando la funzione $f$ è molto variabile su $[a,b]$ o quando
l'intervallo è ampio, perché l'errore dipende dall'ampiezza $h = b - a$
elevata a una potenza alta.
L'idea delle \textbf{formule composite} è:
\begin{enumerate}
  \item suddividere $[a,b]$ in $n$ sottointervalli più piccoli;
  \item applicare la formula semplice su ciascun sottointervallo;
  \item sommare i contributi.
\end{enumerate}

Poiché su ogni sottointervallo il passo locale $h = (b-a)/n$ diventa
piccolo al crescere di $n$, l'errore totale diminuisce rapidamente.

\subsection{Trapezio Composto}

\subsubsection*{Impostazione}

Dati $n+1$ nodi $\{x_0, x_1, \ldots, x_n\}$ equidistanti in $[a,b]$
con passo $h = (b-a)/n$, suddividiamo l'intervallo nei sottointervalli:
\[
  T_k = [x_k, x_{k+1}], \quad k = 0, 1, \ldots, n-1
\]

\noindent In particolare: $T_1 = [x_0,x_1]$, $T_2 = [x_1,x_2]$, \ldots,
$T_n = [x_{n-1},x_n]$.

\subsubsection*{Somma dei trapezi}

Applicando la formula del trapezio~\eqref{eq:trapezio} su ciascun
sottointervallo $T_k$ (larghezza $h$):
\[
  \int_{x_k}^{x_{k+1}} f(x)\,dx \approx \frac{h}{2}\bigl[f(x_k)+f(x_{k+1})\bigr]
\]

Sommando su tutti gli $n$ sottointervalli:
\begin{align*}
  \tilde{I} &\approx
    \frac{h}{2}[f(x_0)+f(x_1)]
  + \frac{h}{2}[f(x_1)+f(x_2)]
  + \cdots
  + \frac{h}{2}[f(x_{n-1})+f(x_n)]
\end{align*}

Raccogliendo i termini: ogni nodo interno $x_i$ (con $1 \le i \le n-1$)
compare due volte (come estremo destro di $T_{i-1}$ e come estremo
sinistro di $T_i$), mentre $x_0$ e $x_n$ compaiono una sola volta.
Quindi:

\begin{equation}
  \boxed{
    \tilde{I} = \frac{h}{2}\!\left(f(x_0) + 2\sum_{i=1}^{n-1} f(x_i) + f(x_n)\right)
  }
  \label{eq:trapezio-comp}
\end{equation}

\subsubsection*{Errore del trapezio composto}

L'errore su ogni sottointervallo $T_k$ è, dalla formula~\eqref{eq:errore-trapezio}:
\[
  R_k = -\frac{h^3}{12}\,f''(\xi_k), \quad \xi_k \in (x_k, x_{k+1})
\]

\noindent Sostituendo $h = (b-a)/n$, si ha $h^3 = (b-a)^3/n^3$, e quindi
l'errore su ogni singolo $T_k$ è:
\[
  R_k = -\frac{(b-a)^3}{12n^3}\,f''(\xi_k)
\]

L'errore totale è la somma su tutti gli $n$ sottointervalli:
\begin{align*}
  R &= \sum_{k=0}^{n-1} R_k
     = \sum_{k=0}^{n-1} \left(-\frac{(b-a)^3}{12n^3}\,f''(\xi_k)\right)
     = -\frac{(b-a)^3}{12n^3}\sum_{k=0}^{n-1} f''(\xi_k)
\end{align*}

\noindent Per il teorema del valor medio applicato alla somma di $n$
valori di $f''$ in punti dell'intervallo, esiste $\xi \in (a,b)$ tale che:
\[
  \sum_{k=0}^{n-1} f''(\xi_k) \approx n\,f''(\xi)
\]

\noindent Di conseguenza:
\[
  R \approx -\frac{(b-a)^3}{12n^3} \cdot n\,f''(\xi) = -\frac{(b-a)^3}{12n^2}\,f''(\xi)
\]

\begin{equation}
  \boxed{R \approx -\frac{(b-a)^3}{12n^2}\,f''(\xi)}
  \label{eq:errore-trap-comp}
\end{equation}

\noindent L'errore decresce come $1/n^2$: raddoppiando il numero di nodi
$n$, l'errore si riduce di un fattore $4$.

\subsection{Cavalieri-Simpson Composto}

\subsubsection*{Impostazione: perché serve un numero pari di nodi}

La formula di Cavalieri-Simpson opera su gruppi di 3 nodi consecutivi.
Per poter suddividere i nodi in tali gruppi senza sovrapposizioni né
buchi, è necessario che $n$ sia \textbf{pari}: in questo caso si hanno
$m = n/2$ sottointervalli (ognuno costituito da 3 nodi).

I sottointervalli ``composti'' sono:
\[
  T_k = [x_{2k},\, x_{2k+1},\, x_{2k+2}]
      = [x_{2k},x_{2k+1}] \cup [x_{2k+1},x_{2k+2}],
      \quad k = 0, 1, \ldots, \frac{n}{2}-1
\]

\noindent dove $x_{2k+1}$ è il punto medio del macro-sottointervallo
$[x_{2k},x_{2k+2}]$, che ha ampiezza $2h$.

\begin{intuizione}
Ogni $T_k$ è un ``macro-sottointervallo'' di ampiezza $2h$, suddiviso
in due intervalli elementari di ampiezza $h$, con il nodo $x_{2k+1}$
come punto medio. Su ogni $T_k$ si applica Cavalieri-Simpson a 3 punti.
\end{intuizione}

\subsubsection*{Esempi al variare di $n$ (con $m = n/2$)}

\begin{center}
\begin{tabular}{ccl}
\toprule
$n$ & $m$ & Sottointervalli $T_k$ \\
\midrule
$2$ & $1$ & $T_0 = [x_0, x_1, x_2]$ \\
$4$ & $2$ & $T_0 = [x_0,x_1,x_2]$, \quad $T_1 = [x_2,x_3,x_4]$ \\
$6$ & $3$ & $T_0 = [x_0,x_1,x_2]$, \quad $T_1 = [x_2,x_3,x_4]$,
            \quad $T_2 = [x_4,x_5,x_6]$ \\
\bottomrule
\end{tabular}
\end{center}

\subsubsection*{Formula generale}

Applicando Cavalieri-Simpson~\eqref{eq:simpson} su ogni $T_k$ e
sommando su $k = 0, 1, \ldots, m-1$:
\[
  \tilde{I}
  \approx \frac{h}{3}\!\left[
    f(x_0)
    + 4\sum_{i=0}^{m-1} f(x_{2i+1})
    + 2\sum_{i=1}^{m-1} f(x_{2i})
    + f(x_n)
  \right]
\]

\noindent Come nel caso del trapezio composto, i nodi agli estremi dei
macro-sottointervalli ($x_0, x_2, x_4, \ldots, x_{n-2}, x_n$) compaiono
più volte: quelli interni ($x_2, x_4, \ldots, x_{n-2}$) vengono sommati
con coefficiente~$2$, mentre i nodi centrali di ogni $T_k$
($x_1, x_3, \ldots, x_{n-1}$) vengono sommati con coefficiente~$4$
(il peso del punto medio nella formula semplice).

Sostituendo $h = (b-a)/n$:

\begin{equation}
  \boxed{
    \tilde{I}
    = \frac{b-a}{6}\!\left[
      f(x_0)
      + 4\sum_{i=0}^{\frac{n}{2}-1} f(x_{2i+1})
      + 2\sum_{i=1}^{\frac{n}{2}-1} f(x_{2i})
      + f(x_n)
    \right]
  }
  \label{eq:simpson-comp}
\end{equation}

\noindent Il 4 che moltiplica la prima sommatoria è ereditato
direttamente dalla formula originale di Cavalieri-Simpson, dove il
punto medio ha peso quattro volte maggiore degli estremi.

\subsubsection*{Errore del Cavalieri-Simpson composto}

L'errore su ogni singolo sottointervallo $T_k$ (di ampiezza $2h$, con
passo locale $h$) è:
\[
  R_k = -\frac{h^5}{90}\,f^{(4)}(\xi_k), \quad \xi_k \in T_k
\]

\noindent dove $f^{(4)}$ indica la derivata quarta di $f$.

L'errore totale si ottiene sommando su tutti gli $m = n/2$
sottointervalli. Sostituendo $h = (b-a)/n$, si ha $h^5 = (b-a)^5/n^5$:
\begin{align*}
  R &= \sum_{k=0}^{m-1} R_k
     \approx m \cdot \left(-\frac{h^5}{90}\right) f^{(4)}(\xi) \\
  &= \frac{n}{2} \cdot \left(-\frac{(b-a)^5}{90\,n^5}\right) f^{(4)}(\xi)
   = -\frac{n\,(b-a)^5}{2 \cdot 90\,n^5}\,f^{(4)}(\xi)
\end{align*}

\begin{equation}
  \boxed{R \approx -\frac{(b-a)^5}{180\,n^4}\,f^{(4)}(\xi)}
  \label{eq:errore-simpson-comp}
\end{equation}

\noindent L'errore decresce come $1/n^4$: raddoppiando i nodi, l'errore
si riduce di un fattore $16$.

\begin{osservazione}[Confronto tra gli errori]
Mettendo a confronto le due formule composite, dove $h = (b-a)/n$:
\begin{center}
\begin{tabular}{lll}
\toprule
Formula & Ordine dell'errore & Fattore di riduzione per $n \to 2n$ \\
\midrule
Trapezio composto       & $\mathcal{O}(h^2) = \mathcal{O}(1/n^2)$ & $\times 1/4$ \\
Cavalieri-Simpson comp. & $\mathcal{O}(h^4) = \mathcal{O}(1/n^4)$ & $\times 1/16$ \\
\bottomrule
\end{tabular}
\end{center}

\noindent La formula di Cavalieri-Simpson ha grado di precisione $p = 3$:
è esatta per tutti i polinomi di grado $\le 3$, pur avendo solo 3 nodi
(che garantirebbero, in linea di principio, solo grado~2). Questo guadagno
di un grado extra è dovuto alla simmetria della formula rispetto al punto
medio dell'intervallo, che fa sì che gli errori di ordine dispari si
cancellino esattamente.
\end{osservazione}

\end{document}
